%%
% 第三章：实验过程与结果
%%

\chapter{实验过程与结果}
\label{chap:experiment}

\section{实验环境}

\begin{table}[htbp]
    \centering
    \caption{实验环境配置}
    \label{tab:env}
    \begin{tabular}{ll}
        \toprule
        \textbf{项目} & \textbf{配置} \\
        \midrule
        操作系统 & Windows 11 Pro \\
        Python版本 & 3.10+ \\
        依赖 & 无第三方依赖（使用标准库） \\
        开发工具 & VS Code / PyCharm \\
        \bottomrule
    \end{tabular}
\end{table}

\section{测试目标程序}

\subsection{程序简介}

本实验使用\texttt{agent\_pageindex.py}作为测试目标程序。该程序是一个基于LangChain的文档问答代理，具有以下特点：

\begin{itemize}
    \item 多层条件判断（参数验证、类型检查）
    \item 循环结构（消息处理、工具调用遍历）
    \item 函数调用（包含嵌套调用）
    \item 异常处理（API密钥检查）
\end{itemize}

\subsection{程序结构分析}

目标程序共130行代码，包含：

\begin{itemize}
    \item \texttt{\_extract\_text()}：文本提取辅助函数
    \item \texttt{\_print\_verbose\_stream()}：详细输出流处理函数
    \item \texttt{main()}：主函数，包含参数解析和Agent执行逻辑
\end{itemize}

\section{实验过程}

\subsection{运行测试}

使用PathTracer命令行工具对目标程序进行追踪：

\begin{lstlisting}[language=bash, caption=测试命令]
# 生成HTML格式报告
python -m pathtracer.cli -f html \
    -o coverage_report.html \
    agent_pageindex.py -- --query "test"

# 生成JSON格式报告
python -m pathtracer.cli -f json \
    -o coverage_report.json \
    agent_pageindex.py -- --query "test"
\end{lstlisting}

\subsection{测试用例设计}

为了验证工具功能，设计了以下测试用例：

\begin{table}[htbp]
    \centering
    \caption{测试用例}
    \label{tab:testcase}
    \begin{tabular}{lll}
        \toprule
        \textbf{编号} & \textbf{测试内容} & \textbf{参数} \\
        \midrule
        TC1 & 基本查询功能 & --query "test" \\
        TC2 & 详细输出模式 & --query "test" --verbose \\
        TC3 & 自定义参数 & --query "test" --doc-top-k 5 \\
        \bottomrule
    \end{tabular}
\end{table}

\section{实验结果}

\subsection{覆盖率统计}

执行TC1测试用例后的覆盖率统计结果：

\begin{table}[htbp]
    \centering
    \caption{覆盖率统计结果}
    \label{tab:coverage}
    \begin{tabular}{ll}
        \toprule
        \textbf{指标} & \textbf{数值} \\
        \midrule
        总分支数 & 22 \\
        已执行分支数 & 7 \\
        覆盖率 & 31.8\% \\
        \bottomrule
    \end{tabular}
\end{table}

\subsection{分支类型分布}

\begin{table}[htbp]
    \centering
    \caption{分支类型分布}
    \label{tab:branch-type}
    \begin{tabular}{lcc}
        \toprule
        \textbf{分支类型} & \textbf{总数} & \textbf{已执行} \\
        \midrule
        if语句 & 17 & 5 \\
        for循环 & 5 & 2 \\
        while循环 & 0 & 0 \\
        except块 & 0 & 0 \\
        \bottomrule
    \end{tabular}
\end{table}

\subsection{已执行分支详情}

\begin{table}[htbp]
    \centering
    \caption{已执行分支列表}
    \label{tab:executed}
    \begin{tabular}{cccl}
        \toprule
        \textbf{行号} & \textbf{分支类型} & \textbf{结束行} & \textbf{代码内容} \\
        \midrule
        25 & if & 26 & payload is None检查 \\
        27 & if & 28 & isinstance(payload, str)检查 \\
        88 & if & 89 & API密钥检查 \\
        116 & if & 118 & verbose模式检查 \\
        121 & if & 125 & result类型检查 \\
        123 & if & 125 & messages存在检查 \\
        129 & if & 130 & \_\_main\_\_入口检查 \\
        \bottomrule
    \end{tabular}
\end{table}

\subsection{函数调用树}

实验过程中记录的函数调用关系：

\begin{lstlisting}[caption=函数调用树]
main()
  load_dotenv()
  os.getenv("DEEPSEEK_API_KEY")
  init_chat_model()
  create_agent()
  agent.invoke()
  _extract_text()
  print()
\end{lstlisting}

\section{结果分析}

\subsection{覆盖率分析}

覆盖率为31.8\%，主要原因：

\begin{enumerate}
    \item \textbf{verbose模式未触发}：\texttt{\_print\_verbose\_stream}函数（40-75行）未被调用，因为测试时未使用\texttt{--verbose}参数
    \item \textbf{异常分支未覆盖}：API密钥存在性检查的异常分支（89行）未触发
    \item \textbf{复杂消息处理未执行}：\texttt{\_extract\_text}函数中的列表处理分支未完全覆盖
    \item \textbf{边缘条件}：部分边缘条件的处理逻辑未在正常流程中执行
\end{enumerate}

\subsection{改进建议}

为提高覆盖率，可以：

\begin{enumerate}
    \item 添加\texttt{--verbose}参数测试
    \item 模拟无API密钥场景
    \item 测试不同类型的消息内容
    \item 编写单元测试覆盖各分支
\end{enumerate}

\section{本章小结}

本章详细记录了实验过程和结果。通过对\texttt{agent\_pageindex.py}的测试，验证了PathTracer工具的有效性。工具成功识别了22个控制流分支，追踪了程序执行路径，计算出31.8\%的覆盖率，并生成了详细的HTML和JSON格式报告。
